% Options for packages loaded elsewhere
\PassOptionsToPackage{unicode}{hyperref}
\PassOptionsToPackage{hyphens}{url}
\documentclass[
]{article}
\usepackage{xcolor}
\usepackage[margin=1in]{geometry}
\usepackage{amsmath,amssymb}
\setcounter{secnumdepth}{-\maxdimen} % remove section numbering
\usepackage{iftex}
\ifPDFTeX
  \usepackage[T1]{fontenc}
  \usepackage[utf8]{inputenc}
  \usepackage{textcomp} % provide euro and other symbols
\else % if luatex or xetex
  \usepackage{unicode-math} % this also loads fontspec
  \defaultfontfeatures{Scale=MatchLowercase}
  \defaultfontfeatures[\rmfamily]{Ligatures=TeX,Scale=1}
\fi
\usepackage{lmodern}
\ifPDFTeX\else
  % xetex/luatex font selection
\fi
% Use upquote if available, for straight quotes in verbatim environments
\IfFileExists{upquote.sty}{\usepackage{upquote}}{}
\IfFileExists{microtype.sty}{% use microtype if available
  \usepackage[]{microtype}
  \UseMicrotypeSet[protrusion]{basicmath} % disable protrusion for tt fonts
}{}
\makeatletter
\@ifundefined{KOMAClassName}{% if non-KOMA class
  \IfFileExists{parskip.sty}{%
    \usepackage{parskip}
  }{% else
    \setlength{\parindent}{0pt}
    \setlength{\parskip}{6pt plus 2pt minus 1pt}}
}{% if KOMA class
  \KOMAoptions{parskip=half}}
\makeatother
\usepackage{color}
\usepackage{fancyvrb}
\newcommand{\VerbBar}{|}
\newcommand{\VERB}{\Verb[commandchars=\\\{\}]}
\DefineVerbatimEnvironment{Highlighting}{Verbatim}{commandchars=\\\{\}}
% Add ',fontsize=\small' for more characters per line
\usepackage{framed}
\definecolor{shadecolor}{RGB}{248,248,248}
\newenvironment{Shaded}{\begin{snugshade}}{\end{snugshade}}
\newcommand{\AlertTok}[1]{\textcolor[rgb]{0.94,0.16,0.16}{#1}}
\newcommand{\AnnotationTok}[1]{\textcolor[rgb]{0.56,0.35,0.01}{\textbf{\textit{#1}}}}
\newcommand{\AttributeTok}[1]{\textcolor[rgb]{0.13,0.29,0.53}{#1}}
\newcommand{\BaseNTok}[1]{\textcolor[rgb]{0.00,0.00,0.81}{#1}}
\newcommand{\BuiltInTok}[1]{#1}
\newcommand{\CharTok}[1]{\textcolor[rgb]{0.31,0.60,0.02}{#1}}
\newcommand{\CommentTok}[1]{\textcolor[rgb]{0.56,0.35,0.01}{\textit{#1}}}
\newcommand{\CommentVarTok}[1]{\textcolor[rgb]{0.56,0.35,0.01}{\textbf{\textit{#1}}}}
\newcommand{\ConstantTok}[1]{\textcolor[rgb]{0.56,0.35,0.01}{#1}}
\newcommand{\ControlFlowTok}[1]{\textcolor[rgb]{0.13,0.29,0.53}{\textbf{#1}}}
\newcommand{\DataTypeTok}[1]{\textcolor[rgb]{0.13,0.29,0.53}{#1}}
\newcommand{\DecValTok}[1]{\textcolor[rgb]{0.00,0.00,0.81}{#1}}
\newcommand{\DocumentationTok}[1]{\textcolor[rgb]{0.56,0.35,0.01}{\textbf{\textit{#1}}}}
\newcommand{\ErrorTok}[1]{\textcolor[rgb]{0.64,0.00,0.00}{\textbf{#1}}}
\newcommand{\ExtensionTok}[1]{#1}
\newcommand{\FloatTok}[1]{\textcolor[rgb]{0.00,0.00,0.81}{#1}}
\newcommand{\FunctionTok}[1]{\textcolor[rgb]{0.13,0.29,0.53}{\textbf{#1}}}
\newcommand{\ImportTok}[1]{#1}
\newcommand{\InformationTok}[1]{\textcolor[rgb]{0.56,0.35,0.01}{\textbf{\textit{#1}}}}
\newcommand{\KeywordTok}[1]{\textcolor[rgb]{0.13,0.29,0.53}{\textbf{#1}}}
\newcommand{\NormalTok}[1]{#1}
\newcommand{\OperatorTok}[1]{\textcolor[rgb]{0.81,0.36,0.00}{\textbf{#1}}}
\newcommand{\OtherTok}[1]{\textcolor[rgb]{0.56,0.35,0.01}{#1}}
\newcommand{\PreprocessorTok}[1]{\textcolor[rgb]{0.56,0.35,0.01}{\textit{#1}}}
\newcommand{\RegionMarkerTok}[1]{#1}
\newcommand{\SpecialCharTok}[1]{\textcolor[rgb]{0.81,0.36,0.00}{\textbf{#1}}}
\newcommand{\SpecialStringTok}[1]{\textcolor[rgb]{0.31,0.60,0.02}{#1}}
\newcommand{\StringTok}[1]{\textcolor[rgb]{0.31,0.60,0.02}{#1}}
\newcommand{\VariableTok}[1]{\textcolor[rgb]{0.00,0.00,0.00}{#1}}
\newcommand{\VerbatimStringTok}[1]{\textcolor[rgb]{0.31,0.60,0.02}{#1}}
\newcommand{\WarningTok}[1]{\textcolor[rgb]{0.56,0.35,0.01}{\textbf{\textit{#1}}}}
\usepackage{graphicx}
\makeatletter
\newsavebox\pandoc@box
\newcommand*\pandocbounded[1]{% scales image to fit in text height/width
  \sbox\pandoc@box{#1}%
  \Gscale@div\@tempa{\textheight}{\dimexpr\ht\pandoc@box+\dp\pandoc@box\relax}%
  \Gscale@div\@tempb{\linewidth}{\wd\pandoc@box}%
  \ifdim\@tempb\p@<\@tempa\p@\let\@tempa\@tempb\fi% select the smaller of both
  \ifdim\@tempa\p@<\p@\scalebox{\@tempa}{\usebox\pandoc@box}%
  \else\usebox{\pandoc@box}%
  \fi%
}
% Set default figure placement to htbp
\def\fps@figure{htbp}
\makeatother
\setlength{\emergencystretch}{3em} % prevent overfull lines
\providecommand{\tightlist}{%
  \setlength{\itemsep}{0pt}\setlength{\parskip}{0pt}}
\usepackage{bookmark}
\IfFileExists{xurl.sty}{\usepackage{xurl}}{} % add URL line breaks if available
\urlstyle{same}
\hypersetup{
  pdftitle={Two proportions{,} chi-square{,} risk and odds},
  hidelinks,
  pdfcreator={LaTeX via pandoc}}

\title{Two proportions, chi-square, risk and odds}
\author{}
\date{\vspace{-2.5em}}

\begin{document}
\maketitle

\textbf{Inference labs} use the five-step template: Hypothesis →
Visualise → Assumptions → Conduct → Conclude.

\subsection{Learning objectives}\label{learning-objectives}

\begin{itemize}
\tightlist
\item
  Compare two proportions with \texttt{prop.test()} and
  \texttt{fisher.test()} and say when each is preferred.
\item
  Compute risk ratio, odds ratio, and risk difference with confidence
  intervals.
\item
  Run a chi-square goodness-of-fit test against a named distribution.
\end{itemize}

\subsection{Prerequisites}\label{prerequisites}

Labs 2.4, 3.4.

\subsection{Background}\label{background}

Two-proportion tests ask whether the rate of an event differs between
two groups. In a 2x2 table, two exposure groups (say, treatment and
control) each produce a fraction of events. The chi-square test compares
the observed counts to those expected under independence; Fisher's exact
test computes the probability of a table as extreme or more extreme
directly from the hypergeometric distribution and is preferred when
expected cell counts are small.

The effect can be expressed in several ways. The risk difference is the
difference of the two proportions. The risk ratio is their ratio --- the
event rate in the treated group, relative to control. The odds ratio is
the ratio of the odds. In a case-control study the risk ratio is not
directly estimable; the odds ratio is. In a cohort study both are
estimable and the risk ratio is usually more intuitive.

Goodness-of-fit tests compare a vector of observed counts to a vector of
expected proportions under a named distribution --- whether a die is
fair, whether a set of genotype frequencies matches Hardy-Weinberg,
whether arrivals over an hour are Poisson. The test statistic is the
same chi-square; the interpretation is ``is the distribution consistent
with a named model?''

\subsection{Setup}\label{setup}

\begin{Shaded}
\begin{Highlighting}[]
\FunctionTok{library}\NormalTok{(tidyverse)}
\FunctionTok{set.seed}\NormalTok{(}\DecValTok{42}\NormalTok{)}
\FunctionTok{theme\_set}\NormalTok{(}\FunctionTok{theme\_minimal}\NormalTok{(}\AttributeTok{base\_size =} \DecValTok{12}\NormalTok{))}
\end{Highlighting}
\end{Shaded}

\subsection{1. Hypothesis}\label{hypothesis}

Two-by-two: \texttt{treatment} vs \texttt{control}, outcome
\texttt{event} vs \texttt{none}.

\begin{itemize}
\tightlist
\item
  H0: event rate is independent of treatment.
\item
  H1: rates differ.
\item
  α = 0.05.
\end{itemize}

\subsection{2. Visualise}\label{visualise}

Simulate: 200 per arm, event rate 0.25 in control, 0.15 in treatment.

\begin{Shaded}
\begin{Highlighting}[]
\NormalTok{sim }\OtherTok{\textless{}{-}} \FunctionTok{tibble}\NormalTok{(}
  \AttributeTok{arm =} \FunctionTok{rep}\NormalTok{(}\FunctionTok{c}\NormalTok{(}\StringTok{"control"}\NormalTok{, }\StringTok{"treatment"}\NormalTok{), }\AttributeTok{each =}\NormalTok{ n\_per),}
  \AttributeTok{event =} \FunctionTok{c}\NormalTok{(}\FunctionTok{rbinom}\NormalTok{(n\_per, }\DecValTok{1}\NormalTok{, }\FloatTok{0.25}\NormalTok{),}
            \FunctionTok{rbinom}\NormalTok{(n\_per, }\DecValTok{1}\NormalTok{, }\FloatTok{0.15}\NormalTok{))}
\NormalTok{)}
\NormalTok{tab }\OtherTok{\textless{}{-}} \FunctionTok{table}\NormalTok{(}\AttributeTok{arm =}\NormalTok{ sim}\SpecialCharTok{$}\NormalTok{arm, }\AttributeTok{event =}\NormalTok{ sim}\SpecialCharTok{$}\NormalTok{event)}
\NormalTok{tab}
\end{Highlighting}
\end{Shaded}

\begin{Shaded}
\begin{Highlighting}[]
\FunctionTok{as\_tibble}\NormalTok{(tab) }\SpecialCharTok{|\textgreater{}}
  \FunctionTok{mutate}\NormalTok{(}\AttributeTok{arm =} \FunctionTok{factor}\NormalTok{(arm, }\AttributeTok{levels =} \FunctionTok{c}\NormalTok{(}\StringTok{"control"}\NormalTok{, }\StringTok{"treatment"}\NormalTok{)),}
         \AttributeTok{event =} \FunctionTok{recode}\NormalTok{(event, }\StringTok{\textasciigrave{}}\AttributeTok{0}\StringTok{\textasciigrave{}} \OtherTok{=} \StringTok{"no event"}\NormalTok{, }\StringTok{\textasciigrave{}}\AttributeTok{1}\StringTok{\textasciigrave{}} \OtherTok{=} \StringTok{"event"}\NormalTok{)) }\SpecialCharTok{|\textgreater{}}
  \FunctionTok{ggplot}\NormalTok{(}\FunctionTok{aes}\NormalTok{(arm, n, }\AttributeTok{fill =}\NormalTok{ event)) }\SpecialCharTok{+}
  \FunctionTok{geom\_col}\NormalTok{(}\AttributeTok{position =} \StringTok{"fill"}\NormalTok{, }\AttributeTok{alpha =} \FloatTok{0.7}\NormalTok{) }\SpecialCharTok{+}
  \FunctionTok{labs}\NormalTok{(}\AttributeTok{x =} \ConstantTok{NULL}\NormalTok{, }\AttributeTok{y =} \StringTok{"Proportion"}\NormalTok{, }\AttributeTok{fill =} \ConstantTok{NULL}\NormalTok{)}
\end{Highlighting}
\end{Shaded}

\subsection{3. Assumptions}\label{assumptions}

Chi-square requires expected cell counts ≥ 5. Fisher has no such
constraint and is exact at any size. Random assignment or representative
sampling is assumed. Independence of observations is essential.

\begin{Shaded}
\begin{Highlighting}[]

\end{Highlighting}
\end{Shaded}

All expected counts large; chi-square is appropriate.

\subsection{4. Conduct}\label{conduct}

\begin{Shaded}
\begin{Highlighting}[]
\NormalTok{ft }\OtherTok{\textless{}{-}} \FunctionTok{fisher.test}\NormalTok{(tab)}
\NormalTok{pt }\OtherTok{\textless{}{-}} \FunctionTok{prop.test}\NormalTok{(}\AttributeTok{x =}\NormalTok{ tab[, }\StringTok{"1"}\NormalTok{], }\AttributeTok{n =} \FunctionTok{rowSums}\NormalTok{(tab), }\AttributeTok{correct =} \ConstantTok{FALSE}\NormalTok{)}
\FunctionTok{list}\NormalTok{(}\AttributeTok{chisq =}\NormalTok{ ct, }\AttributeTok{fisher =}\NormalTok{ ft, }\AttributeTok{prop =}\NormalTok{ pt)}
\end{Highlighting}
\end{Shaded}

\subsubsection{Effect sizes: RD, RR, OR}\label{effect-sizes-rd-rr-or}

\begin{Shaded}
\begin{Highlighting}[]
\NormalTok{p\_trt  }\OtherTok{\textless{}{-}}\NormalTok{ tab[}\StringTok{"treatment"}\NormalTok{, }\StringTok{"1"}\NormalTok{] }\SpecialCharTok{/} \FunctionTok{sum}\NormalTok{(tab[}\StringTok{"treatment"}\NormalTok{, ])}

\NormalTok{rd }\OtherTok{\textless{}{-}}\NormalTok{ p\_trt }\SpecialCharTok{{-}}\NormalTok{ p\_ctrl}

\NormalTok{rr }\OtherTok{\textless{}{-}}\NormalTok{ p\_trt }\SpecialCharTok{/}\NormalTok{ p\_ctrl}
\CommentTok{\# Log{-}based Wald CI for RR}
\NormalTok{se\_log\_rr }\OtherTok{\textless{}{-}} \FunctionTok{sqrt}\NormalTok{(}
\NormalTok{  (}\DecValTok{1} \SpecialCharTok{{-}}\NormalTok{ p\_trt)  }\SpecialCharTok{/}\NormalTok{ (p\_trt  }\SpecialCharTok{*} \FunctionTok{sum}\NormalTok{(tab[}\StringTok{"treatment"}\NormalTok{, ])) }\SpecialCharTok{+}
\NormalTok{  (}\DecValTok{1} \SpecialCharTok{{-}}\NormalTok{ p\_ctrl) }\SpecialCharTok{/}\NormalTok{ (p\_ctrl }\SpecialCharTok{*} \FunctionTok{sum}\NormalTok{(tab[}\StringTok{"control"}\NormalTok{,   ]))}
\NormalTok{)}
\NormalTok{ci\_rr }\OtherTok{\textless{}{-}} \FunctionTok{exp}\NormalTok{(}\FunctionTok{log}\NormalTok{(rr) }\SpecialCharTok{+} \FunctionTok{c}\NormalTok{(}\SpecialCharTok{{-}}\DecValTok{1}\NormalTok{, }\DecValTok{1}\NormalTok{) }\SpecialCharTok{*} \FunctionTok{qnorm}\NormalTok{(}\FloatTok{0.975}\NormalTok{) }\SpecialCharTok{*}\NormalTok{ se\_log\_rr)}

\NormalTok{or }\OtherTok{\textless{}{-}}\NormalTok{ (tab[}\StringTok{"treatment"}\NormalTok{, }\StringTok{"1"}\NormalTok{] }\SpecialCharTok{*}\NormalTok{ tab[}\StringTok{"control"}\NormalTok{, }\StringTok{"0"}\NormalTok{]) }\SpecialCharTok{/}
\NormalTok{      (tab[}\StringTok{"treatment"}\NormalTok{, }\StringTok{"0"}\NormalTok{] }\SpecialCharTok{*}\NormalTok{ tab[}\StringTok{"control"}\NormalTok{, }\StringTok{"1"}\NormalTok{])}
\NormalTok{se\_log\_or }\OtherTok{\textless{}{-}} \FunctionTok{sqrt}\NormalTok{(}\FunctionTok{sum}\NormalTok{(}\DecValTok{1} \SpecialCharTok{/}\NormalTok{ tab))}
\NormalTok{ci\_or }\OtherTok{\textless{}{-}} \FunctionTok{exp}\NormalTok{(}\FunctionTok{log}\NormalTok{(or) }\SpecialCharTok{+} \FunctionTok{c}\NormalTok{(}\SpecialCharTok{{-}}\DecValTok{1}\NormalTok{, }\DecValTok{1}\NormalTok{) }\SpecialCharTok{*} \FunctionTok{qnorm}\NormalTok{(}\FloatTok{0.975}\NormalTok{) }\SpecialCharTok{*}\NormalTok{ se\_log\_or)}

\FunctionTok{tibble}\NormalTok{(}\AttributeTok{effect =} \FunctionTok{c}\NormalTok{(}\StringTok{"RD"}\NormalTok{, }\StringTok{"RR"}\NormalTok{, }\StringTok{"OR"}\NormalTok{),}
       \AttributeTok{est =} \FunctionTok{c}\NormalTok{(rd, rr, or),}
       \AttributeTok{ci\_low =} \FunctionTok{c}\NormalTok{(pt}\SpecialCharTok{$}\NormalTok{conf.int[}\DecValTok{1}\NormalTok{], ci\_rr[}\DecValTok{1}\NormalTok{], ci\_or[}\DecValTok{1}\NormalTok{]),}
       \AttributeTok{ci\_high =} \FunctionTok{c}\NormalTok{(pt}\SpecialCharTok{$}\NormalTok{conf.int[}\DecValTok{2}\NormalTok{], ci\_rr[}\DecValTok{2}\NormalTok{], ci\_or[}\DecValTok{2}\NormalTok{]))}
\end{Highlighting}
\end{Shaded}

\subsubsection{Goodness-of-fit example}\label{goodness-of-fit-example}

Test whether 600 simulated genotype counts match Hardy-Weinberg
proportions at allele frequency 0.3.

\begin{Shaded}
\begin{Highlighting}[]
\NormalTok{pAA }\OtherTok{\textless{}{-}} \FloatTok{0.7}\SpecialCharTok{\^{}}\DecValTok{2}\NormalTok{; paa }\OtherTok{\textless{}{-}} \FloatTok{0.3}\SpecialCharTok{\^{}}\DecValTok{2}\NormalTok{; pAa }\OtherTok{\textless{}{-}} \DecValTok{2} \SpecialCharTok{*} \FloatTok{0.7} \SpecialCharTok{*} \FloatTok{0.3}
\NormalTok{exp\_probs }\OtherTok{\textless{}{-}} \FunctionTok{c}\NormalTok{(}\AttributeTok{AA =}\NormalTok{ pAA, }\AttributeTok{Aa =}\NormalTok{ pAa, }\AttributeTok{aa =}\NormalTok{ paa)}
\NormalTok{gof }\OtherTok{\textless{}{-}} \FunctionTok{chisq.test}\NormalTok{(obs, }\AttributeTok{p =}\NormalTok{ exp\_probs)}
\NormalTok{gof}
\end{Highlighting}
\end{Shaded}

\subsection{5. Concluding statement}\label{concluding-statement}

\begin{quote}
Event rates were \texttt{round(p\_ctrl,\ 3)} in control and
\texttt{round(p\_trt,\ 3)} in treatment (RD \texttt{round(rd,\ 3)}; RR
\texttt{round(rr,\ 2)}, 95\% CI \texttt{round(ci\_rr{[}1{]},\ 2)} to
\texttt{round(ci\_rr{[}2{]},\ 2)}; OR \texttt{round(or,\ 2)}, 95\% CI
\texttt{round(ci\_or{[}1{]},\ 2)} to \texttt{round(ci\_or{[}2{]},\ 2)}).
Chi-square test χ² = \texttt{round(ct\$statistic,\ 2)}, df =
\texttt{ct\$parameter}, \emph{p} = \texttt{signif(ct\$p.value,\ 2)}. The
Hardy-Weinberg goodness-of-fit test on simulated genotypes gave χ² =
\texttt{round(gof\$statistic,\ 2)} (df = \texttt{gof\$parameter},
\emph{p} = \texttt{signif(gof\$p.value,\ 2)}).
\end{quote}

Report RR or OR alongside the \emph{p}-value. In a cohort design prefer
RR, which is the quantity clinicians translate directly to practice.

\subsection{Common pitfalls}\label{common-pitfalls}

\begin{itemize}
\tightlist
\item
  Reporting an OR when a RR is meaningful and more interpretable.
\item
  Using the continuity correction (the default in \texttt{prop.test} and
  \texttt{chisq.test}) by habit; it is conservative.
\item
  Running a chi-square when expected counts are small; use Fisher.
\item
  Treating independent observations as if repeated measures weren't.
\end{itemize}

\subsection{Further reading}\label{further-reading}

\begin{itemize}
\tightlist
\item
  Altman DG. \emph{Practical Statistics for Medical Research}, chapter
  10--11.
\item
  Agresti A. \emph{Categorical Data Analysis}.
\end{itemize}

\subsection{Session info}\label{session-info}

\begin{Shaded}
\begin{Highlighting}[]

\end{Highlighting}
\end{Shaded}

\subsection{See also --- chapter index}\label{see-also-chapter-index}

\begin{itemize}
\tightlist
\item
  \href{../index.Rmd}{Methods in this chapter}
\item
  \href{../../../ch00-find-your-method.Rmd}{Find your method (Chapter
  0)}
\end{itemize}

\end{document}
